\section{Valuation rings}
\label{section-valuation-rings}

\noindent
Here are some definitions.

\begin{definition}
\label{definition-valuation-ring}
Valuation rings.
\begin{enumerate}
\item Let $K$ be a field. Let $A$, $B$ be local rings contained
in $K$. We say that $B$ {\it dominates} $A$ if $A \subset B$
and $\mathfrak m_A = A \cap \mathfrak m_B$.
\item Let $A$ be a ring. We say $A$ is a {\it valuation ring}
if $A$ is a local domain and if $A$ is maximal
for the relation of domination among local rings contained in
the fraction field of $A$.
\item Let $A$ be a valuation ring with fraction field $K$.
If $R \subset K$ is a subring of $K$, then we say $A$
is {\it centered} on $R$ if $R \subset A$.
\end{enumerate}
\end{definition}

\noindent
With this definition a field is a valuation ring.

\begin{lemma}
\label{lemma-dominate}
Let $K$ be a field. Let $A \subset K$ be a local subring.
Then there exists a valuation ring with fraction field $K$
dominating $A$.
\end{lemma}

\begin{proof}
We consider the collection of local subrings
of $K$ as a partially ordered set using the relation of domination.
Suppose that $\{A_i\}_{i \in I}$ is a totally ordered
collection of local subrings of $K$. Then $B = \bigcup A_i$
is a local subring which dominates all of the $A_i$.
Hence by Zorn's Lemma, it suffices to show that if $A \subset K$
is a local ring whose fraction field is not $K$, then there
exists a local ring $B \subset K$, $B \not = A$ dominating $A$.

\medskip\noindent
Pick $t \in K$ which is not in the fraction field of $A$.
If $t$ is transcendental over $A$, then $A[t] \subset K$
and hence $A[t]_{(t, \mathfrak m)} \subset K$ is a local ring
distinct from $A$ dominating $A$. Suppose $t$ is algebraic over $A$.
Then for some nonzero $a \in A$ the element $at$ is integral over $A$.
In this case the subring $A' \subset K$ generated by $A$ and
$ta$ is finite over $A$.
By Lemma \ref{lemma-integral-overring-surjective} there exists
a prime ideal $\mathfrak m' \subset A'$ lying over
$\mathfrak m$. Then $A'_{\mathfrak m'}$ dominates
$A$. If $A = A'_{\mathfrak m'}$, then $t$
is in the fraction field of $A$ which we assumed not to be the case.
Thus $A \not = A'_{\mathfrak m'}$ as desired.
\end{proof}

\begin{lemma}
\label{lemma-valuation-ring-normal}
Let $A$ be a valuation ring.
Then $A$ is a normal domain.
\end{lemma}

\begin{proof}
Suppose $x$ is in the field of fractions of $A$ and integral over $A$.
Let $A'$ denote the subring of $K$ generated by $A$ and $x$.
Since $A\subset A'$ is an integral extension, we see by
Lemma \ref{lemma-integral-overring-surjective} that there is
a prime ideal $\mathfrak m' \subset A'$ lying over
$\mathfrak m$. Then $A'_{\mathfrak m'}$ dominates
$A$. Since $A$ is a valuation ring we conclude that $A=A'_{\mathfrak m'}$
and therefore that $x\in A$.
\end{proof}


\begin{lemma}
\label{lemma-valuation-ring-x-or-x-inverse}
Let $A$ be a valuation ring with maximal ideal $\mathfrak m$ and
fraction field $K$.
Let $x \in K$. Then either $x \in A$ or $x^{-1} \in A$ or both.
\end{lemma}

\begin{proof}
Assume that $x$ is not in $A$.
Let $A'$ denote the subring of $K$ generated by $A$ and $x$.
Since $A$ is a valuation ring we see that there is no prime
of $A'$ lying over $\mathfrak m$. Since $\mathfrak m$ is maximal
we see that $V(\mathfrak m A') = \emptyset$. Then $\mathfrak m A' = A'$
by Lemma \ref{lemma-Zariski-topology}.
Hence we can write
$1 = \sum_{i = 0}^d t_i x^i$ with $t_i \in \mathfrak m$.
This implies that $(1 - t_0) (x^{-1})^d - \sum t_i (x^{-1})^{d - i} = 0$.
In particular we see that $x^{-1}$ is integral over $A$, and hence
$x^{-1} \in A$ by
Lemma \ref{lemma-valuation-ring-normal}.
\end{proof}

\begin{lemma}
\label{lemma-x-or-x-inverse-valuation-ring}
Let $A \subset K$ be a subring of a field $K$ such that for all
$x \in K$ either $x \in A$ or $x^{-1} \in A$ or both.
Then $A$ is a valuation ring with fraction field $K$.
\end{lemma}

\begin{proof}
If $A$ is not $K$, then $A$ is not a field and there is a nonzero
maximal ideal $\mathfrak m$.
If $\mathfrak m'$ is a second maximal ideal, then choose $x, y \in A$
with $x \in \mathfrak m$, $y \not \in \mathfrak m$,
$x \not \in \mathfrak m'$, and $y \in \mathfrak m'$.
Then neither $x/y \in A$ nor $y/x \in A$
contradicting the assumption of the lemma. Thus we see that $A$ is
a local ring. Suppose that $A'$ is a local ring contained in $K$ which
dominates $A$. Let $x \in A'$. We have to show that $x \in A$. If not, then
$x^{-1} \in A$, and of course $x^{-1} \in \mathfrak m_A$. But then
$x^{-1} \in \mathfrak m_{A'}$ which contradicts $x \in A'$.
\end{proof}

\begin{lemma}
\label{lemma-colimit-valuation-rings}
\begin{slogan}
Valuation rings are stable under filtered direct limits
\end{slogan}
Let $I$ be a directed set. Let $(A_i, \varphi_{ij})$
be a system of valuation rings over $I$.
Then $A = \colim A_i$ is a valuation ring.
\end{lemma}

\begin{proof}
It is clear that $A$ is a domain. Let $a, b \in A$.
Lemma \ref{lemma-x-or-x-inverse-valuation-ring} tells us we have
to show that either $a | b$ or $b | a$ in $A$. Choose $i$
so large that there exist $a_i, b_i \in A_i$ mapping to $a, b$.
Then Lemma \ref{lemma-valuation-ring-x-or-x-inverse}
applied to $a_i, b_i$ in $A_i$ implies the result for $a, b$ in $A$.
\end{proof}

\begin{lemma}
\label{lemma-valuation-ring-cap-field}
Let $L/K$ be an extension of fields. If $B \subset L$
is a valuation ring, then $A = K \cap B$ is a valuation ring.
\end{lemma}

\begin{proof}
We can replace $L$ by the fraction field $F$ of $B$ and $K$ by
$K \cap F$. Then the lemma follows from a combination of
Lemmas \ref{lemma-valuation-ring-x-or-x-inverse} and
\ref{lemma-x-or-x-inverse-valuation-ring}.
\end{proof}

\begin{lemma}
\label{lemma-valuation-ring-cap-field-finite}
Let $L/K$ be an algebraic extension of fields. If $B \subset L$
is a valuation ring with fraction field $L$ and not a field, then
$A = K \cap B$ is a valuation ring and not a field.
\end{lemma}

\begin{proof}
By Lemma \ref{lemma-valuation-ring-cap-field} the ring $A$ is a valuation
ring. If $A$ is a field, then $A = K$. Then $A = K \subset B$ is an integral
extension, hence there are no proper inclusions among the primes of $B$
(Lemma \ref{lemma-integral-no-inclusion}).
This contradicts the assumption that $B$ is a local domain and not a field.
\end{proof}

\begin{lemma}
\label{lemma-make-valuation-rings}
Let $A$ be a valuation ring. For any prime ideal $\mathfrak p \subset A$ the
quotient $A/\mathfrak p$ is a valuation ring. The same is true for the
localization $A_\mathfrak p$ and in fact any localization of $A$.
\end{lemma}

\begin{proof}
Use the characterization of valuation rings given
in Lemma \ref{lemma-x-or-x-inverse-valuation-ring}.
\end{proof}

\begin{lemma}
\label{lemma-stack-valuation-rings}
Let $A'$ be a valuation ring with residue field $K$.
Let $A$ be a valuation ring with fraction field $K$.
Then
$C = \{\lambda \in A' \mid \lambda \bmod \mathfrak m_{A'} \in A\}$
is a valuation ring.
\end{lemma}

\begin{proof}
Note that $\mathfrak m_{A'} \subset C$ and $C/\mathfrak m_{A'} = A$.
In particular, the fraction field of $C$ is equal to the fraction field
of $A'$. We will use the criterion of
Lemma \ref{lemma-x-or-x-inverse-valuation-ring} to prove the lemma.
Let $x$ be an element of the fraction field of $C$.
By the lemma we may assume $x \in A'$. If $x \in \mathfrak m_{A'}$,
then we see $x \in C$. If not, then $x$ is a unit of $A'$ and we
also have $x^{-1} \in A'$. Hence either $x$ or $x^{-1}$ maps to
an element of $A$ by the lemma again.
\end{proof}

\begin{lemma}
\label{lemma-find-valuation-rings}
Let $A$ be a normal domain with fraction field $K$.
\begin{enumerate}
\item For every $x \in K$, $x \not \in A$ there exists a valuation ring
$A \subset V \subset K$ with fraction field $K$ such that $x \not \in V$.
\item If $A$ is local, we can moreover choose $V$ which dominates $A$.
\end{enumerate}
In other words, $A$ is the intersection of all valuation rings in $K$
containing $A$ and if $A$ is local, then $A$ is the intersection of
all valuation rings in $K$ dominating $A$.
\end{lemma}

\begin{proof}
Suppose $x \in K$, $x \not \in A$. Consider $B = A[x^{-1}]$.
Then $x \not \in B$. Namely, if $x = a_0 + a_1x^{-1} + \ldots + a_d x^{-d}$
then $x^{d + 1} - a_0x^d - \ldots - a_d = 0$ and $x$ is integral
over $A$ in contradiction with the fact that $A$ is normal.
Thus $x^{-1}$ is not a unit in $B$. Thus $V(x^{-1}) \subset \Spec(B)$
is not empty (Lemma \ref{lemma-Zariski-topology}), and we can choose a prime
$\mathfrak p \subset B$ with $x^{-1} \in \mathfrak p$.
Choose a valuation ring $V \subset K$ dominating $B_\mathfrak p$
(Lemma \ref{lemma-dominate}).
Then $x \not \in V$ as $x^{-1} \in \mathfrak m_V$.

\medskip\noindent
If $A$ is local, then we claim that $x^{-1} B + \mathfrak m_A B \not = B$.
Namely, if $1 = (a_0 + a_1x^{-1} + \ldots + a_d x^{-d})x^{-1} +
a'_0 + \ldots + a'_d x^{-d}$ with $a_i \in A$ and $a'_i \in \mathfrak m_A$,
then we'd get
$$
(1 - a'_0) x^{d + 1} - (a_0 + a'_1) x^d - \ldots - a_d = 0
$$
Since $a'_0 \in \mathfrak m_A$ we see that $1 - a'_0$ is a unit in $A$
and we conclude that $x$ would be integral over $A$, a contradiction as
before. Then choose the prime $\mathfrak p \supset x^{-1} B + \mathfrak m_A B$
we find $V$ dominating $A$.
\end{proof}

\noindent
An {\it totally ordered abelian group} is a pair $(\Gamma, \geq)$
consisting of an abelian group $\Gamma$ endowed with a total
ordering $\geq$ such that $\gamma \geq \gamma' \Rightarrow
\gamma + \gamma'' \geq \gamma' + \gamma''$ for all
$\gamma, \gamma', \gamma'' \in \Gamma$.

\begin{lemma}
\label{lemma-valuation-group}
Let $A$ be a valuation ring with field of fractions $K$.
Set $\Gamma = K^*/A^*$ (with group law written additively).
For $\gamma, \gamma' \in \Gamma$
define $\gamma \geq \gamma'$ if and only if
$\gamma - \gamma'$ is in the image of $A - \{0\} \to \Gamma$.
Then $(\Gamma, \geq)$ is a totally ordered abelian group.
\end{lemma}

\begin{proof}
Omitted, but follows easily from
Lemma \ref{lemma-valuation-ring-x-or-x-inverse}.
Note that in case $A = K$ we obtain the zero group $\Gamma = \{0\}$
endowed with its unique total ordering.
\end{proof}

\begin{definition}
\label{definition-value-group}
Let $A$ be a valuation ring.
\begin{enumerate}
\item The totally ordered abelian group $(\Gamma, \geq)$ of
Lemma \ref{lemma-valuation-group} is called the
{\it value group} of the valuation ring $A$.
\item The map $v : A - \{0\} \to \Gamma$ and also $v : K^* \to \Gamma$ is
called the {\it valuation} associated to $A$.
\item The valuation ring $A$ is called a {\it discrete valuation ring}
if $\Gamma \cong \mathbf{Z}$.
\end{enumerate}
\end{definition}

\noindent
Note that if $\Gamma \cong \mathbf{Z}$ then there is a unique such
isomorphism such that $1 \geq 0$. If the isomorphism is chosen in this
way, then the ordering becomes the usual ordering of the integers.

\begin{lemma}
\label{lemma-properties-valuation}
Let $A$ be a valuation ring. The valuation $v : A -\{0\} \to \Gamma_{\geq 0}$
has the following properties:
\begin{enumerate}
\item $v(a) = 0 \Leftrightarrow a \in A^*$,
\item $v(ab) = v(a) + v(b)$,
\item $v(a + b) \geq \min(v(a), v(b))$ provided $a + b \not = 0$.
\end{enumerate}
\end{lemma}

\begin{proof}
Omitted.
\end{proof}

\begin{lemma}
\label{lemma-characterize-valuation-ring}
Let $A$ be a ring. The following are equivalent
\begin{enumerate}
\item $A$ is a valuation ring,
\item $A$ is a local domain and every finitely generated
ideal of $A$ is principal.
\end{enumerate}
\end{lemma}

\begin{proof}
Assume $A$ is a valuation ring and let $f_1, \ldots, f_n \in A$.
Choose $i$ such that $v(f_i)$ is minimal among $v(f_j)$.
Then $(f_i) = (f_1, \ldots, f_n)$. Conversely, assume $A$ is
a local domain and every finitely generated ideal of $A$ is principal.
Pick $f, g \in A$ and write $(f, g) = (h)$. Then $f = ah$ and $g = bh$
and $h = cf + dg$ for some $a, b, c, d \in A$. Thus $ac + bd = 1$
and we see that either $a$ or $b$ is a unit, i.e., either
$g/f$ or $f/g$ is an element of $A$. This shows $A$ is a valuation ring
by Lemma \ref{lemma-x-or-x-inverse-valuation-ring}.
\end{proof}

\begin{lemma}
\label{lemma-valuation-valuation-ring}
Let $(\Gamma, \geq)$ be a totally ordered abelian group.
Let $K$ be a field. Let $v : K^* \to \Gamma$ be a homomorphism
of abelian groups such that $v(a + b) \geq \min(v(a), v(b))$ for
$a, b \in K$ with $a, b, a + b$ not zero. Then
$$
A =
\{
x \in K \mid x = 0 \text{ or } v(x) \geq 0
\}
$$
is a valuation ring with value group $\Im(v) \subset \Gamma$,
with maximal ideal
$$
\mathfrak m =
\{
x \in K \mid x = 0 \text{ or } v(x) > 0
\}
$$
and with group of units
$$
A^* =
\{
x \in K^* \mid v(x) = 0
\}.
$$
\end{lemma}

\begin{proof}
Omitted.
\end{proof}

\noindent
Let $(\Gamma, \geq)$ be a totally ordered abelian group.
An {\it ideal of $\Gamma$} is a subset $I \subset \Gamma$ such
that all elements of $I$ are $\geq 0$ and $\gamma \in I$,
$\gamma' \geq \gamma$ implies $\gamma' \in I$. We say that such
an ideal is {\it prime} if $0 \not \in I$ and if
$\gamma + \gamma' \in I, \gamma, \gamma' \geq 0
\Rightarrow \gamma \in I \text{ or } \gamma' \in I$.

\begin{lemma}
\label{lemma-ideals-valuation-ring}
Let $A$ be a valuation ring.
Ideals in $A$ correspond $1 - 1$ with ideals of $\Gamma$.
This bijection is inclusion preserving, and maps prime
ideals to prime ideals.
\end{lemma}

\begin{proof}
Omitted.
\end{proof}

\begin{lemma}
\label{lemma-valuation-ring-Noetherian-discrete}
A valuation ring is Noetherian if and only if it is
a discrete valuation ring or a field.
\end{lemma}

\begin{proof}
Suppose $A$ is a discrete valuation ring
with valuation $v : A \setminus \{0\} \to \mathbf{Z}$
normalized so that $\Im(v) = \mathbf{Z}_{\geq 0}$.
By Lemma \ref{lemma-ideals-valuation-ring} the ideals of $A$ are the subsets
$I_n = \{0\} \cup v^{-1}(\mathbf{Z}_{\geq n})$. It is clear
that any element $x \in A$ with $v(x) = n$ generates $I_n$.
Hence $A$ is a PID so certainly Noetherian.

\medskip\noindent
Suppose $A$ is a Noetherian valuation ring with value group $\Gamma$.
By Lemma \ref{lemma-ideals-valuation-ring} we see the ascending chain
condition holds for ideals in $\Gamma$.
We may assume $A$ is not a field, i.e., there is a $\gamma \in \Gamma$
with $\gamma > 0$. Applying the ascending chain condition to the subsets
$\gamma + \Gamma_{\geq 0}$ with $\gamma > 0$ we see
there exists a smallest element $\gamma_0$ which is bigger than $0$.
Let $\gamma \in \Gamma$ be an element $\gamma > 0$. Consider the sequence
of elements $\gamma$, $\gamma - \gamma_0$, $\gamma - 2\gamma_0$,
etc. By the ascending chain condition these cannot all be $> 0$.
Let $\gamma - n \gamma_0$ be the last one $\geq 0$. By minimality
of $\gamma_0$ we see that $0 = \gamma - n \gamma_0$. Hence $\Gamma$
is a cyclic group as desired.
\end{proof}























